\song{Babička Mary}
\artist{voskovecwerich}

\large% the song does not fit on one page at \Large once the block gaps cannot be squeezed

\vers{1}{
\chord{Am}Štěchovická laguna když dřímá v zadumaném stínu Kordiller,\\
\chord{Dm}pirát zkrva\chord{Am}venou šerpu ždímá, \chord{H\7}šerif si láduje revol\chord{E}ver.\\
\chord{Am}Pikovická rýžoviště zlata čeří se v příboji Sázavy,\\
\chord{Dm}ale zato \chord{Am}krčmářova chata \chord{Dm}křepčí rykem \chord{E\7}chlapské zába\chord{Am}vy.\\
\chord{G\7}Když tu náhle, co se \chord{C}děje, \chord{G\7}divný šelest houštím \chord{C}spěje,\\
\chord{G\7}plch, skunk, \chord{C}vše utíká \chord{F}po stráni \chord{F\kk\dm}od Mední\chord{E}ka.\\
\chord{Am}Krčmář zhasne, kovbojové ztichnou, pirát zděšen tvář si zakryje,\\
\chord{Dm}rudé squaw se \chord{Am}chvějí, a pak vzdychnou: \chord{H\7}\uv{Blíží se k nám postrach préri\chord{E}e.} \chords{G\7}
}

\refrain{
\chord{C}Mary, babička \chord{D\7}Mary, dva kolťá\chord{G\7}ky za pasem, nad hlavou \chord{C}točí lasem.\\
Stoletá Mary, babička \chord{D\7}Mary, ta zkrotí \chord{G\7}křepce hřebce, ať chce, či ne\chord{C}chce. \chord{E\7}
}

\vers{2}{
Žádné zuby, z jelenice sukně, ale zato tvrdé bicepsy,\\
Mary má vždy slivovici v putně, Toma Mixe\footnote{Tom Mix (1880–1940) byl hvězdou amerických filmových westernů a symbol neohroženého kovboje a hrdiny Divokého západu.} strčí do kapsy.\\
Klika cvakla, v krčmě dveře letí, a babička vchází do dveří,\\
\uv{Pintu ginu, lumpové prokletí!} bezzubou dásní zaláteří.\\
Vypiju to jen ve stoje, jdu do volebního boje,\\
zřím zas město drahý, jedu volit do Prahy.\\
Dopila a aby se neřeklo, putykáře změní v mrtvolu,\\
za zády má štěchovické peklo, s šlajsnou\footnote{Vodácký slangový výraz pro propust (z němčiny \textit{Schleuse} = plavební komora, zdymadlo nebo propust), tedy upravenou část jezu určenou ke sjíždění lodí nebo pro odtok vody.} svatojánských\footnote{Svatojánské proudy: Bývalý divoký kaňon Vltavy mezi Slapy a Štěchovicemi, který byl kolébkou českého trampingu a vodáctví, než ho v polovině 20. století zatopila Vltavská kaskáda.} atolů\footnote{Atol je typ ostrova v tropických mořích, který se skládá z kruhovitého korálového útesu, jenž obklopuje lagunu. Zde jde o vtipné přirovnání vltavského údolí k exotickému mořskému ráji.}.
}

\refrain{
Mary, babička Mary, pádluje bez námahy, po proudu až do Prahy.\\
Stoletá Mary, babička Mary jde do volebního boje za kovboje.
}

\vers{3}{
Ledva v Praze kotvu vyhodila, pro babičku nastal hrozný čas,\\
neboť každá strana hned tvrdila, že jí náleží babiččin hlas.\\
Malá stejně jako velká strana psala, že bude mít o hlas víc,\\
že ta druhá strana je nahraná, oni že maj hlas ze Štěchovic.\\
Stoletý věk prý nevadí, na předáka je to mládí,\\
ze všech nejvíce, volala ji Polnice.\\
Tak babičku, pro kterou vždy byla válka s lidojedy legrace,\\
tu babička za pár dní zabila volební agitace.
}

\refrain{
Mary, bojovná Mary, už nesedává  v sedle, ve volbách byla vedle.\\
Stoletou Mary, babičku Mary volbama zabitou vzal k sobě Manitou.
}
